\documentclass[12pt,a4paper]{article}
\usepackage{graphicx}
\usepackage{amsmath}
\usepackage{amssymb}
\usepackage{amsthm}
\usepackage{bm}
\usepackage{hyperref}
\usepackage{geometry}
\geometry{margin=2.5cm}
\emergencystretch=3em

\newtheorem{theorem}{Theorem}
\newtheorem{postulate}{Postulate}
\newtheorem{definition}{Definition}

\title{\textbf{CDFD Runtime Paper 01: Axioms, State Variables, and Claim Boundaries}}
\author{Steve Bico Mujjabi, MD\\
\small Independent Researcher\\
\small ORCID: \href{https://orcid.org/0009-0001-0556-5516}{0009-0001-0556-5516}}
\date{May 2026}

\begin{document}
\maketitle

\begin{abstract}
\noindent \textbf{Executive synthesis.} This paper defines the public axiomatic
surface of Constraint-Driven Flux Dynamics (CDFD) as implemented in
\texttt{cdfd\_runtime}. The framework represents a system through active flow
$\Phi$, adaptive constraint $C$, surface responsiveness $S$, retained structural
memory $M_s$, and the operating ratio
$\Psi_s=(\Phi/C)S M_s$. These variables are not presented here as empirical
proof of a universal physical ontology. They are a computational grammar: a
state language that can be stress-tested, falsified, and reused across the
runtime, CDFL models, domain adapters, validation tools, and local discovery
experiments.
\end{abstract}

\paragraph{Greek-symbol convention.}
Unless a manuscript states a tighter equation-local meaning, Greek symbols are local CDFD variables or coefficients, not universal constants. Here $\Phi$ denotes driving flux, $\Psi_s$ denotes the adaptive operating ratio, $\Omega$ denotes a domain or state space, and $\Lambda$ denotes the Life Number where used. The coefficients $\alpha$, $\beta$, and $\gamma$ denote accumulation or reinforcement, relaxation or decay, and diffusion, coupling, or redistribution. The symbols $\delta$ and $\epsilon$ denote perturbation or tolerance scales, while $\eta$, $\lambda$, $\mu$, $\nu$, $\rho$, $\sigma$, $\tau$, and $\phi$ denote local efficiency, rate, baseline, frequency, density or correlation, noise or variance, time-scale, and phase or field terms. Subscripts restrict any coefficient to the named pathway, tissue, domain, or experiment.


\begin{figure}[ht]
\centering
\fbox{\begin{minipage}{0.92\textwidth}
\centering
\textbf{Runtime evidence path}\\[0.4em]
Prior CDFD mathematics $\longrightarrow$ CDFL/runtime state $\longrightarrow$ bounded output $\longrightarrow$ falsification target.
\end{minipage}}
\caption{Conceptual schematic for the runtime papers. The engine translates the mathematical frame from the prior CDFD papers into executable CDFL/runtime diagnostics; candidate outputs remain tied to explicit tests before any stronger claim is made.}
\end{figure}

\section{Prior CDFD Basis}
This manuscript does not rederive the mathematical core of CDFD. It uses the
Mujjabi Adaptive Operating Ratio and the constrained-vacuum notation introduced
in Part I \cite{MujjabiI2026}, the torus-knot hierarchy and boundary-mode work
from the Part I sequence \cite{MujjabiVI2026,MujjabiIX2026}, the public balance
law developed in the Part I capstone \cite{MujjabiX2026}, the Part I transport
tests \cite{MujjabiXII2026}, and the Life Number and tri-regime biology bridge
from Part II \cite{MujjabiPartII2026}. The runtime papers therefore act as an
execution layer: they keep the mathematics connected to commands, outputs,
falsification tests, and domain-specific limits.


\section{Purpose and Scope}
The CDFD Runtime begins from a pragmatic observation: many systems fail when
drive exceeds the capacity of the medium that must carry it. A thermal channel,
a metabolic pathway, a city road network, a financial clearing system, and a
prebiotic mineral pore are materially different. Yet each can be represented,
at a coarse level, as a relation between throughput and constraint.

This paper states the shared variables. Later papers derive dynamics, describe
the runtime, explain CDFL, and connect the theory to local experiments in
\texttt{cdfd\_runtime/experiments}. The claim discipline is deliberately strict:
runtime experiments are candidate model diagnostics and falsification targets
unless independent empirical validation is provided.

\section{Core State Variables}
\begin{definition}[Flow field $\Phi$]
The flow field $\Phi(\bm{x},t)$ is the active drive or transported quantity of a
system. Depending on context, $\Phi$ may represent heat flux, substrate flux,
capital velocity, information flow, pressure drive, or field intensity.
\end{definition}

\begin{definition}[Constraint field $C$]
The constraint field $C(\bm{x},t)$ is the local resistance, bottleneck,
capacity limit, or structural load that opposes or shapes $\Phi$. In CDFD,
$C$ is not fixed by default. It may adapt, relax, diffuse, stiffen, or fail.
\end{definition}

\begin{definition}[Surface responsiveness $S$]
The surface field $S(\bm{x},t)$ describes how readily a boundary, interface, or
medium changes its routing behavior under load. A high $S$ can improve routing
when regulated, but can amplify instability when recovery is weak.
\end{definition}

\begin{definition}[Structural memory $M_s$]
The memory field $M_s(\bm{x},t)$ records retained stiffening after overload.
In the engine this is the hysteresis channel: a system can remain altered after
the original pulse, injury, or shock has passed.
\end{definition}

\begin{definition}[Operating ratio $\Psi_s$]
The adaptive operating ratio is
\begin{equation}
    \Psi_s(\bm{x},t)=\frac{\Phi(\bm{x},t)}{C(\bm{x},t)}S(\bm{x},t)M_s(\bm{x},t).
\end{equation}
Low $\Psi_s$ indicates under-driving or over-constraint. Near-unity $\Psi_s$
marks an actively regulated operating band. High $\Psi_s$ indicates overload,
with interpretation dependent on the modeled domain.
\end{definition}

\section{Axioms}
\begin{postulate}[Flow requires a medium]
Every modeled flow is carried by a medium with finite effective constraint.
The limiting behavior $C \to 0$ is therefore a boundary condition, not an
ordinary engineering state.
\end{postulate}

\begin{postulate}[Constraint responds to stress]
Sustained or rapidly changing $\Phi$ can alter $C$, $S$, and $M_s$. This is the
reason the runtime stores fields rather than isolated scalar labels.
\end{postulate}

\begin{postulate}[Regime transitions are measurable]
The useful test of CDFD is not whether CDFD language can be fitted to a system
after the fact. The useful test is whether changes in $\Phi$, $C$, $S$, and
$M_s$ predict measurable transitions before they occur.
\end{postulate}

\section{Runtime Connection}
These variables are implemented in \texttt{engine/state.py}. The physics update
lives primarily in \texttt{engine/physics.py}; shared finite checks, JSON-safe
summaries, the scalar operating ratio, the Life Number, and bounded diagnostic
helpers live in \texttt{runtime/diagnostics.py}. The command surface is exposed
through \texttt{cdfd.py} (\texttt{info}, \texttt{domains}, \texttt{demo},
\texttt{diagnostics}, \texttt{validate}, \texttt{run}, \texttt{export},
\texttt{auth}), while CDFL models provide a readable way to define systems and
rules. \texttt{cdfd.py diagnostics} writes
\texttt{experiments/outputs/part\_ii\_runtime\_diagnostics.json} with guarded
Part II source-mix and endpoint status rows.

\section{Experiment and Discovery Hooks}
The experiments directory provides candidate diagnostics tied to these axioms:
\begin{itemize}
    \item \texttt{experiments/notebooks/discovery\_vacuum\_hysteresis.py}
    tests whether a high-flux pulse leaves a localized $M_s$ residual.
    \item \texttt{experiments/notebooks/discovery\_adaptive\_surface.py}
    stress-tests when $S$ amplifies rather than stabilizes overload.
    \item \texttt{experiments/notebooks/discovery\_ool\_phase.py}
    explores the Life Number boundary $\Lambda \approx 1$ in a prebiotic toy
    parameter space.
    \item \texttt{experiments/reports/FINAL\_RELEASE\_EXPERIMENT\_SELECTION.md}
    records which outputs are suitable for public discussion and which must
    remain only simulation logs.
\end{itemize}

\section{Claim Status}
The axioms are a computational foundation and a falsifiable modeling proposal.
They are not by themselves a claim that every physical, biological, or social
system has been empirically validated under CDFD. A domain mapping fails if it
does not produce reproducible regime thresholds, if its $\Psi_s$ proxy adds no
predictive information beyond ordinary variables, or if independent data reject
its proposed flow-constraint mapping.

\section{Conclusion}
Paper 01 establishes the vocabulary: $\Phi$, $C$, $S$, $M_s$, and $\Psi_s$.
The rest of the runtime series asks whether this vocabulary can become useful
science: executable, falsifiable, reproducible, and disciplined about the
difference between a model diagnostic and an empirical discovery.
\begin{thebibliography}{9}
\bibitem{MujjabiI2026}
Steve Bico Mujjabi, MD. \emph{Vortex Stability in a Constrained Transport Vacuum and the Origin of the Fine-Structure Constant}. CDFD Part I Paper I, preprint, 2026.

\bibitem{MujjabiVI2026}
Steve Bico Mujjabi, MD. \emph{The Universal Torus Knot Hierarchy: $Z_n$ Symmetry, the Cinquefoil Family, and the General Structure of CDFT Particle Families}. CDFD Part I Paper VI, preprint, 2026.

\bibitem{MujjabiIX2026}
Steve Bico Mujjabi, MD. \emph{Even-$n$ Torus Knots in CDFT: The Exclusion of $T(2,2)$, the Extension to $T(2,2m)$, and the Anti-Phase Mass Scaling Law}. CDFD Part I Paper IX, preprint, 2026.

\bibitem{MujjabiX2026}
Steve Bico Mujjabi, MD. \emph{The Vacuum Equation of State: Deriving Torus Knot Self-Consistency from the Public CDFD Balance Law $\Psi = \Phi/C$}. CDFD Part I Paper X, preprint, 2026.

\bibitem{MujjabiXII2026}
Steve Bico Mujjabi, MD. \emph{Friction, Superconductivity, Plasma States, and Transport Mysteries: Observable Tests of Adaptive Constraint}. CDFD Part I Paper XII, preprint, 2026.

\bibitem{MujjabiPartII2026}
Steve Bico Mujjabi, MD. \emph{CDFD Part II: Origins of Life and Tri-Regime Bioenergetics}. Public release archive, 2026, \url{https://doi.org/10.5281/zenodo.20264779}.
\end{thebibliography}
\end{document}
